% GENERATED FILE. Edit canonical lessons, then run npm run generate:books.
% canonical-source-hash: fnv1a64:6e03b9c86d1f7eea
% canonical-lessons: ZH-C359-guize, ZH-C359-zunshou, ZH-C359-weifan, ZH-C359-xianzhi, ZH-C359-daode

\chapter{Rules}
\label{ch:zh-rules}
\hlchaptermodality{359}

\begin{chapteropening}
\textbf{By the end of this chapter:} \emph{I can talk about keeping and breaking rules, putting limits on things, and what is right.}

Complete the last lesson of chapter 359: I can talk about keeping and breaking rules, putting limits on things, and what is right.
\end{chapteropening}

\section[guīzé]{guīzé — rules of a game; set procedure}

\label{lesson:ZH-C359-guize}



\begin{quote}
Before the new one: say the Mandarin for modest about yourself, then the Mandarin for face (your standing).
\end{quote}

\subsection*{You'll want to know: guīzé}
\emph{guīzé} — \textquotedblleft{}rules of a game; set procedure\textquotedblright{}.

\begin{quote}
\emph{zhège bǐsài de guīzé hěn jiǎndān} — The rules of this game are simple.
\end{quote}

Its characters come later; for now it is a word for the ear.

\begin{cousinweb}
\emph{guīzé} is built from \emph{guī} and \emph{zé}, both inherited from Middle Chinese, \emph{kjwie} and \emph{tsok} (Guangyun, in Baxter's transcription), which glosses the first 'round' and the second 'a rule'.
\end{cousinweb}

\begin{grammarlens}[title={What you've built}]
One more word for giving reasons.
\end{grammarlens}

\subsection*{Your turn}
\begin{itemize}
\raggedright
  \item \emph{Say it:} \emph{guīzé}
  \item \emph{Say it:} \emph{guīzé}, once more
  \item \emph{Say it:} \emph{zhège bǐsài de guīzé hěn jiǎndān} — The rules of this game are simple.
  \item \emph{Recall:} say the Mandarin for modest about yourself, then the Mandarin for face (your standing), then say \emph{guīzé} again
\end{itemize}

\begin{tcolorbox}[breakable,colback=teal!4,colframe=teal!35!black,title={Before you move on}]
What does guīzé mean? (\textquotedblleft{}Rules of a game; set procedure\textquotedblright{}.) Say it once more.
\end{tcolorbox}
\section[zūnshǒu]{zūnshǒu — to keep to; to abide by}

\label{lesson:ZH-C359-zunshou}



\begin{quote}
Before the new one: say the Mandarin for face (your standing), then the Mandarin for the rules of a game.
\end{quote}

\subsection*{You'll want to know: zūnshǒu}
\emph{zūnshǒu} — \textquotedblleft{}to keep to; to abide by\textquotedblright{}.

\begin{quote}
\emph{wǒmen yīnggāi zūnshǒu guīzé} — We should keep to the rules.
\end{quote}

Its characters come later; for now it is a word for the ear.

\begin{cousinweb}
\emph{zūnshǒu} is built from \emph{zūn} and \emph{shǒu}, both inherited from Middle Chinese, \emph{tswin} and \emph{syuwX} (Guangyun, in Baxter's transcription), which glosses the first 'to follow' and the second 'to guard'.
\end{cousinweb}

\begin{grammarlens}[title={What you've built}]
One more word for giving reasons.
\end{grammarlens}

\subsection*{Your turn}
\begin{itemize}
\raggedright
  \item \emph{Say it:} \emph{zūnshǒu}
  \item \emph{Say it:} \emph{zūnshǒu}, once more
  \item \emph{Say it:} \emph{wǒmen yīnggāi zūnshǒu guīzé} — We should keep to the rules.
  \item \emph{Recall:} say the Mandarin for face (your standing), then the Mandarin for the rules of a game, then say \emph{zūnshǒu} again
\end{itemize}

\begin{tcolorbox}[breakable,colback=teal!4,colframe=teal!35!black,title={Before you move on}]
What does zūnshǒu mean? (\textquotedblleft{}To keep to; to abide by\textquotedblright{}.) Say it once more.
\end{tcolorbox}
\section[wéifǎn]{wéifǎn — to break a rule; to go against the law}

\label{lesson:ZH-C359-weifan}



\begin{quote}
Before the new one: say the Mandarin for the rules of a game, then the Mandarin for to abide by.
\end{quote}

\subsection*{You'll want to know: wéifǎn}
\emph{wéifǎn} — \textquotedblleft{}to break a rule; to go against the law\textquotedblright{}.

\begin{quote}
\emph{tā wéifǎn le xuéxiào de guīdìng} — He broke the school rules.
\end{quote}

Its characters come later; for now it is a word for the ear.

\begin{cousinweb}
\emph{wéifǎn} is built from \emph{wéi} and \emph{fǎn}, both inherited from Middle Chinese, \emph{hjw+j} and \emph{pjonX} (Guangyun, in Baxter's transcription), which glosses the first 'to turn one's back on' and the second 'to turn over'.
\end{cousinweb}

\begin{grammarlens}[title={What you've built}]
One more word for giving reasons.
\end{grammarlens}

\subsection*{Your turn}
\begin{itemize}
\raggedright
  \item \emph{Say it:} \emph{wéifǎn}
  \item \emph{Say it:} \emph{wéifǎn}, once more
  \item \emph{Say it:} \emph{tā wéifǎn le xuéxiào de guīdìng} — He broke the school rules.
  \item \emph{Recall:} say the Mandarin for the rules of a game, then the Mandarin for to abide by, then say \emph{wéifǎn} again
\end{itemize}

\begin{tcolorbox}[breakable,colback=teal!4,colframe=teal!35!black,title={Before you move on}]
What does wéifǎn mean? (\textquotedblleft{}To break a rule; to go against the law\textquotedblright{}.) Say it once more.
\end{tcolorbox}
\section[xiànzhì]{xiànzhì — to put a limit on; to restrict}

\label{lesson:ZH-C359-xianzhi}



\begin{quote}
Before the new one: say the Mandarin for to abide by, then the Mandarin for to break a rule.
\end{quote}

\subsection*{You'll want to know: xiànzhì}
\emph{xiànzhì} — \textquotedblleft{}to put a limit on; to restrict\textquotedblright{}.

\begin{quote}
\emph{wǒmen yīnggāi xiànzhì háizi shàngwǎng de shíjiān} — We should limit how long children spend online.
\end{quote}

Its characters come later; for now it is a word for the ear.

\begin{cousinweb}
\emph{xiànzhì} is built from \emph{xiàn} and \emph{zhì}, both inherited from Middle Chinese, \emph{heanX} and \emph{tsyejH} (Guangyun, in Baxter's transcription), which glosses the first 'a measure' and 'a boundary', and the second 'to restrain'.
\end{cousinweb}

\begin{grammarlens}[title={What you've built}]
One more word for giving reasons.
\end{grammarlens}

\subsection*{Your turn}
\begin{itemize}
\raggedright
  \item \emph{Say it:} \emph{xiànzhì}
  \item \emph{Say it:} \emph{xiànzhì}, once more
  \item \emph{Say it:} \emph{wǒmen yīnggāi xiànzhì háizi shàngwǎng de shíjiān} — We should limit how long children spend online.
  \item \emph{Recall:} say the Mandarin for to abide by, then the Mandarin for to break a rule, then say \emph{xiànzhì} again
\end{itemize}

\begin{tcolorbox}[breakable,colback=teal!4,colframe=teal!35!black,title={Before you move on}]
What does xiànzhì mean? (\textquotedblleft{}To put a limit on; to restrict\textquotedblright{}.) Say it once more.
\end{tcolorbox}
\section[dàodé]{dàodé — morality; ethics}

\label{lesson:ZH-C359-daode}



\begin{quote}
Before the new one: say the Mandarin for to break a rule, then the Mandarin for to put a limit on.
\end{quote}

\subsection*{You'll want to know: dàodé}
\emph{dàodé} — \textquotedblleft{}morality; ethics\textquotedblright{}.

\begin{quote}
\emph{wǒ juéde zhè shì dàodé wèntí} — I think it's a moral question.
\end{quote}

Its characters come later; for now it is a word for the ear.

\begin{cousinweb}
\emph{dàodé} is built from \zh{道} 'a way' and \emph{dé} 'virtue', both inherited from Middle Chinese, \emph{dawX} and \emph{tok} (Guangyun, in Baxter's transcription), which glosses the first 'a principle' and the second 'virtuous conduct'.
\end{cousinweb}

\begin{grammarlens}[title={What you've built}]
One more word for giving reasons.
\end{grammarlens}

\subsection*{Your turn}
\begin{itemize}
\raggedright
  \item \emph{Say it:} \emph{dàodé}
  \item \emph{Say it:} \emph{dàodé}, once more
  \item \emph{Say it:} \emph{wǒ juéde zhè shì dàodé wèntí} — I think it's a moral question.
  \item \emph{Recall:} say the Mandarin for to break a rule, then the Mandarin for to put a limit on, then say \emph{dàodé} again
\end{itemize}

\begin{tcolorbox}[breakable,colback=teal!4,colframe=teal!35!black,title={Before you move on}]
What does dàodé mean? (\textquotedblleft{}Morality; ethics\textquotedblright{}.) Say it once more.
\end{tcolorbox}
